With high energy density, a long cycle life, and a low self-discharge rate, lithium-ion batteries have become one of the most important electrochemical energy storage technologies, widely used across consumer electronics, electric vehicles, stationary energy storage, and other domains\cite{ref1,ref2,ref3,ref4}. Nevertheless, lithium-ion batteries can still pose considerable safety risks during long-term operation. Charge-discharge cycling and complex operating conditions induce aging processes such as thickening of the solid electrolyte interphase (SEI) layer, loss of active lithium, and structural degradation of the electrodes\cite{ref5,ref6}. These processes, in turn, lead to capacity fade and performance degradation and, in severe cases, trigger failures such as short circuits and thermal runaway\cite{ref7,ref8}. Therefore, precise estimation of lithium-ion battery state of health (SOH) is essential for ensuring battery safety, reliability, and operational performance.

Capacity fade is the most direct sign of battery aging, and state of health (SOH) is commonly defined in terms of capacity retention\cite{ref13}. However, accurately measuring a battery's maximum available capacity requires full or nearly full charge-discharge cycles, which are impractical for online monitoring\cite{ref10}. Consequently, indirectly inferring SOH from observable operating signals has become a key approach to online battery health monitoring. Nevertheless, complex operating conditions make it harder to consistently extract degradation information, while the limited computing power and storage space of battery management systems (BMS) constrain model complexity. Together, these factors pose a dual challenge for online SOH estimation: degradation information extraction and the computational efficiency of estimation models.

\subsection{Literature review}

Recent approaches to SOH estimation can generally be classified as model-based or data-driven\cite{ref12}. The former include electrochemical models and equivalent circuit models, while the latter include traditional machine learning and deep learning methods.

Model-based approaches simulate the internal electrochemical mechanisms of batteries using mathematical equations or equivalent circuits\cite{ref19}. Electrochemical models use differential equations to represent electrochemical reactions and degradation processes in batteries. For instance, Li et al.\cite{ref21} developed an SOH estimation framework based on a single-particle model (SPM), incorporating both chemical degradation and mechanical damage. Although these models are physically interpretable and can achieve high estimation accuracy, solving their complex equations requires extensive computational resources. In addition, parameters such as particle radius and diffusion coefficients are difficult to obtain. These limitations hinder the online application of electrochemical models in resource-constrained BMS\cite{ref20,ref22,ref23}. In contrast, equivalent circuit models (ECMs) simplify complex electrochemical processes into circuits composed of resistors, capacitors, and other components to simulate battery charge-discharge dynamics. For example, Chen et al.\cite{ref24} used recursive least squares to identify the parameters of a Thevenin model and estimated SOH based on the relationship between ohmic internal resistance and capacity fade. This simplification reduces computational cost, but it also limits the ability of ECMs to fully capture changes in the internal state of the battery. Compared with electrochemical models, ECMs may achieve lower estimation accuracy. Their accuracy depends on the selected circuit structure and model parameters. Furthermore, ECM parameters are highly sensitive to operating conditions, including temperature and charge-discharge rate, making it difficult for ECMs to maintain robust estimation performance throughout the battery aging process.

Data-driven approaches do not require a detailed electrochemical model or an extensive analysis of internal reactions and aging mechanisms\cite{ref11}. They train models on historical operating data to map extracted health indicators (HIs) to SOH estimates. Earlier studies primarily employed traditional machine learning models to estimate battery states and predict battery lifetime\cite{ref25}. For instance, Fei et al.\cite{ref26} crafted 42 features from charge-discharge data collected over the first 100 cycles. After feature selection, the retained features were fed into six models: elastic net, Gaussian process regression (GPR), support vector machine (SVM), random forest (RF), gradient boosting regression tree (GBRT), and neural network (NN). Their performance in the early prediction of battery lifetime was then compared. Li et al.\cite{ref27} used a random forest to estimate battery capacity online from partial charging voltage--capacity data. On strongly nonlinear, noisy cycling data from online monitoring and historical cycles, traditional machine-learning models often provide limited predictive performance because of their structural limitations.

With the rapid development of deep learning, researchers have increasingly adopted multilayer neural networks. These networks have strong nonlinear representational capacity and can therefore capture complex battery degradation patterns and improve estimation accuracy\cite{ref14}. Common deep learning models include convolutional neural networks (CNNs), recurrent neural networks (RNNs), and Transformers. CNNs extract local features using convolutional kernels and reduce feature dimensionality through pooling operations. For example, Qian et al.\cite{ref29} fed randomly selected charging-curve segments into a one-dimensional CNN for battery capacity estimation. The results showed that the method achieved low capacity estimation errors with randomly selected segments as inputs. Lee et al.\cite{ref30} transformed capacity fade sequences across cycles into two-dimensional images and employed a two-dimensional CNN for SOH estimation.

However, a single convolutional layer in a traditional CNN has a limited local receptive field, and fusing distant features generally requires multiple convolutional layers\cite{ref31}. RNNs establish these long-range connections by passing historical information to subsequent time steps through hidden states. RNN variants, including long short-term memory (LSTM) and gated recurrent unit (GRU) networks, further employ gating mechanisms to selectively retain and update historical information to capture long-term dependencies. Researchers have also combined CNNs with recurrent networks to capture local features while retaining historical information. For example, Xu et al.\cite{ref41} fed CNN-extracted features into an LSTM and used its temporal memory to estimate SOH. Tian et al.\cite{ref42} further incorporated an attention mechanism into a CNN-BiLSTM architecture to develop the CNN-BiLSTM-AM prediction model. Zheng et al.\cite{ref33} employed a CNN-GRU to estimate SOH using segments of voltage, current, and temperature curves obtained during random charging. Nevertheless, the hidden states in these models still require sequential updates across time steps, limiting parallel computation during training and inference.

To address this limitation, Vaswani et al.\cite{ref34} proposed the Transformer, which avoids recurrence and uses self-attention to directly connect different sequence positions, capturing long-range dependencies while supporting parallel computation. For example, Park and Kim\cite{ref38} used a Transformer guided by physical priors to estimate long-term battery SOH, while Chen et al.\cite{ref37} added structures such as dimension-transformation layers to a vision Transformer to adapt it to battery SOH estimation. Researchers have also combined Transformers with convolutional or recurrent networks to enhance prediction accuracy. Gu et al.\cite{ref35} used a CNN to extract local features and a Transformer to capture long-range dependencies. Jia et al.\cite{ref32} combined a BiGRU with a Transformer, and their experiments showed good prediction accuracy, robustness, and generalization on the battery data used. Bai and Wang\cite{ref36} proposed a convolutional Transformer framework that extracts local details from voltage and current signals and uses self-attention to establish global connections among local features. However, many existing methods improve Transformer prediction accuracy by integrating other modules or models, often at the cost of more parameters and higher computational overhead. Meanwhile, the computational bottleneck of standard self-attention remains: it computes pairwise relationships among all sequence positions, with time and memory complexity growing quadratically with sequence length $N$, namely $O(N^2)$\cite{ref39,ref40}. These computational and storage demands make it challenging to use such models online in resource-constrained BMS\cite{ref31}.

In addition to the prediction model, health indicator extraction and selection are crucial for accurate SOH estimation. These indicators are mainly extracted from voltage, current, temperature, time, and internal resistance data recorded during battery cycling, as well as curves derived from these data\cite{ref47,ref49}. Some internal-resistance indicators require impedance spectrum measurements, which are time-consuming and require specialized instruments. Although temperature indicators contain battery degradation information, they are also affected by ambient temperature and charge-discharge rate, limiting their stability across operating conditions\cite{ref64}. Incremental capacity curves reflect changes in voltage plateaus caused by battery aging, but their calculation requires differentiation, which can amplify sampling noise. Smoothing is therefore usually needed to obtain reliable curve features. For example, Wen et al.\cite{ref50} extracted candidate features such as peak values, peak positions, and peak-shape slopes from smoothed incremental capacity curves and used correlation analysis to select the peak value as the input to a backpropagation (BP) neural network.

In contrast, time-based health indicators do not require differentiation and are simpler to extract. Constant current charge time (CCCT) can be obtained from the difference between the times corresponding to the two endpoints of a selected voltage interval. Richardson et al.\cite{ref28} used charging-time features from different voltage intervals of constant-current charging curves to estimate battery capacity, while Lin et al.\cite{ref63} used CCCT as the input to a random forest to estimate battery SOH. Because the correlation between charging-time features and capacity depends on the selected voltage interval, researchers generally need to identify a suitable interval for feature extraction. Tian et al.\cite{ref51} optimized the charging voltage interval, strengthening the correlation between the extracted health indicator and capacity. In selecting intervals for charging-time features, Li et al.\cite{ref31} progressively narrowed the voltage window and evaluated candidates using Pearson correlation coefficient (PCC) values from multiple cells in the same group, aiming to extract charging-time features highly correlated with SOH from shorter charging segments.

Different types of operating data and the curves derived from them reflect different aspects of battery degradation. To use this information more fully, some studies construct candidate health indicators from multiple operating signals and further optimize model inputs. For example, Dai et al.\cite{ref48} extracted health indicators from voltage, current, temperature, incremental capacity, and differential thermal voltammetry curves and calculated statistical features such as the per-cycle means and medians of these indicators. They then compared estimation errors for different combinations of statistical features to determine the inputs for SOH estimation. Lin et al.\cite{ref55} constructed multiple features from electrical, thermodynamic, and electrochemical perspectives, used principal component analysis to reduce dimensionality, and optimized the number of retained dimensions with simulated annealing, thereby reducing redundancy and improving SOH estimation performance.

A single health indicator contains relatively limited degradation information, whereas multi-source features can describe complementary aspects of battery degradation. However, adding more candidate indicators does not necessarily improve estimation accuracy. Indicators with weak representational ability, inconsistent performance across cells, or redundant information may limit the effective use of multi-source information and increase the model's learning burden. Therefore, a health indicator selection algorithm should be developed to jointly evaluate candidates' correlations with SOH, performance across cells, and redundancy among indicators. It should retain indicators with stable representational ability and reduce inputs weakly correlated with SOH or containing redundant information, helping the model learn battery degradation patterns more effectively.

A comparison of existing SOH estimation methods is summarized in \Cref{tab:soh-method-comparison}. The table compares the source signals, HI-selection strategies, model architectures, and whether computational efficiency is evaluated. Given differences in battery chemistries and operating conditions, the ability of selected indicators to represent degradation consistently across cells warrants further attention. In addition to estimation performance, model complexity is directly related to practical deployment. In recent years, a notable trend in deep model design has been to increase network depth, expand model size, or introduce more complex structures for feature interaction to enhance representational capacity. These designs help models better capture complex degradation patterns but usually come with higher parameter counts, computational costs, and storage requirements. Therefore, SOH estimation models need to become more compact and efficient while retaining effective representational capacity to meet the real-time estimation requirements of resource-constrained BMS.

\begin{table}
\TableStyle
\caption{Comparison of existing SOH estimation methods.}
\label{tab:soh-method-comparison}
\begin{tabular*}{\linewidth}{@{\extracolsep{\fill}}>{\raggedright\arraybackslash}p{40pt}>{\centering\arraybackslash}p{31pt}>{\centering\arraybackslash}p{31pt}>{\centering\arraybackslash}p{48pt}>{\raggedright\arraybackslash}p{110pt}>{\raggedright\arraybackslash}p{111pt}>{\centering\arraybackslash}p{50pt}@{}}
\toprule
Reference & \multicolumn{3}{c}{Source data considered} & HI selection method (author-designed) & Estimation model (original design) & Efficiency considered \\
\cmidrule(lr){2-4}
 & Current & Voltage & Temperature & & & \\
\midrule
\cite{ref78} & \ding{51} & \ding{51} & \ding{55} & PCC (\ding{55}) & PINN (\ding{51}) & \ding{55} \\
\cite{ref31} & \ding{55} & \ding{51} & \ding{55} & PCC (\ding{51}) & BMSFormer (\ding{51}) & \ding{51} \\
\cite{ref32} & \ding{55} & \ding{51} & \ding{51} & PCC + SCC (\ding{55}) & BiGRU-Transformer (\ding{55}) & \ding{55} \\
\cite{ref37} & \ding{51} & \ding{51} & \ding{55} & PCC (\ding{55}) & Transformer (\ding{55}) & \ding{55} \\
\cite{ref79} & \ding{51} & \ding{51} & \ding{55} & FOI monotonicity analysis (\ding{51}) & Linear regression (\ding{55}) & \ding{55} \\
\cite{ref27} & \ding{55} & \ding{51} & \ding{55} & LC (\ding{55}) & Random Forest Regression (\ding{55}) & \ding{55} \\
\cite{ref35} & \ding{51} & \ding{51} & \ding{51} & PCC (\ding{55}) & CNN-Transformer (\ding{55}) & \ding{51} \\
\cite{ref55} & \ding{51} & \ding{51} & \ding{51} & SA–PCA (\ding{51}) & CNN (\ding{55}) & \ding{51} \\
\cite{ref80} & \ding{55} & \ding{51} & \ding{55} & PCC (\ding{55}) & BiLSTM-Transformer (\ding{51}) & \ding{55} \\
\cite{ref81} & \ding{55} & \ding{51} & \ding{55} & Interval selection + feature fusion (\ding{51}) & GWO-LSSVM-AdaBoost (\ding{51}) & \ding{55} \\
\cite{ref82} & \ding{55} & \ding{51} & \ding{51} & PCC + input comparison (\ding{55}) & AST-LSTM (\ding{51}) & \ding{55} \\
\cite{ref83} & \ding{51} & \ding{51} & \ding{51} & PCC (\ding{55}) & LSTM-LAM (\ding{51}) & \ding{55} \\
\cite{ref84} & \ding{51} & \ding{51} & \ding{51} & PCC + SCC (\ding{55}) & BAS-ELM (\ding{51}) & \ding{55} \\
\cite{ref63} & \ding{55} & \ding{51} & \ding{55} & CCCT window search (\ding{51}) & Random Forest Regression (\ding{55}) & \ding{51} \\
\cite{ref48} & \ding{51} & \ding{51} & \ding{51} & Statistical feature combination optimization (\ding{51}) & NGO-Dualkernel-GPR (\ding{51}) & \ding{55} \\
\cite{ref51} & \ding{51} & \ding{51} & \ding{55} & QPSO interval optimization (\ding{51}) & LSTM-BP (\ding{55}) & \ding{55} \\
This study & \ding{51} & \ding{51} & \ding{51} & Multi-source health indicator extraction and optimization algorithm (\ding{51}) & MS-AgentNet (\ding{51}) & \ding{51} \\
\bottomrule
\end{tabular*}
\end{table}

\FloatBarrier

\subsection{Challenges and contributions}

Three main challenges facing existing methods are summarized as follows.

(1) \textbf{Cross-cell stability of health indicators.} Raw data collected from multiple sensors contain diverse features related to degradation. The relationship between a candidate health indicator and SOH may vary across cells, and an indicator that performs well on one cell may not retain the same representational ability on others. PCC and the Spearman correlation coefficient (SCC) evaluate linear and monotonic relationships, respectively; using only one makes it difficult to assess both aspects. Moreover, even indicators strongly associated with SOH may contain redundant information. Model input selection therefore needs to consider both correlations across cells and redundancy among indicators to reduce its dependence on the performance of a single cell.

(2) \textbf{Fusion of local and long-term degradation information.} As cycling progresses, battery capacity generally declines, with fluctuations between neighboring cycles and local capacity recovery. Local features reflect short-term changes, but capturing long-term capacity fade trends also requires connections across cycles. When aggregating information across cycles, local degradation details must be retained so that short-term changes are not weakened. Therefore, the model needs not only to extract degradation features at different time scales but also to make local variations and long-term trends complement each other to improve SOH estimation accuracy.

(3) \textbf{Computational complexity limitations.} Many existing deep learning models improve SOH estimation accuracy by increasing network depth or combining different network structures. These approaches increase the number of model parameters and computational operations. Recurrent structures require sequential computation over time steps, while standard self-attention computes pairwise relationships among all sequence positions, with time and memory complexity growing quadratically with sequence length. These costs increase the operating burden on resource-constrained BMS and limit the deployment of online SOH estimation models.

To address these challenges, this paper proposes a lithium-ion battery SOH estimation framework built around MS-AgentNet. The framework contains two components: group-level health-indicator calibration and selection, and a lightweight network that models local variations and longer-term trends with low computational overhead. The main contributions are as follows.

(1) \textbf{A multi-source health indicator extraction and optimization algorithm is proposed.} The algorithm first extracts multiple types of candidate health indicators from charge-discharge data and their derived curves. MS-CCCT then calibrates the voltage window for the charging-time feature. The algorithm uses correlation results from the feature-development cells to evaluate each candidate's association with SOH and its redundancy with other indicators. It then applies PCC and SCC thresholds and redundancy constraints to select the indicators. The same selection rules determine a final HI subset for each dataset, and the selected indicators retain strong linear and monotonic correlations with SOH on other cells in that dataset.

(2) \textbf{A lightweight local-global network, MS-AgentNet, is constructed.} To represent both long-term battery degradation trends and short-term capacity fluctuations, small-kernel depthwise separable convolutions are combined with ReLU$^2$ agent attention to form a Local-Global Fusion Attention (LGFA) module, integrating local feature extraction and global context modeling with low parameter overhead. The model further introduces large-kernel depthwise separable convolutions, which work with the small-kernel convolutions to extract degradation features at different time scales. With a fixed number of agents and feature dimension, its attention computation reduces the complexity with respect to sequence length from $O(N^2)$ for standard self-attention to $O(N)$. These designs aim to balance SOH estimation accuracy with computational and storage overhead.

(3) \textbf{Comprehensive validation is conducted on four public datasets.} The method is evaluated on four public battery datasets with different chemistries, capacities, and cycling protocols. These experiments assess estimation accuracy, cross-cell generalization, ablation behavior, and computational complexity. Cross-dataset transfer experiments further examine the model's transfer performance across different data domains. The results show that the proposed model maintains high SOH estimation accuracy with low computational and storage overhead.
